% --- Archon physics formalization source begin ---
% archon:physics
% archon:covers PhyXMiniProblems/problem_phyx_mini_0286.lean
% archon:source-report reports/phyx_mini/problem_phyx_mini_0286.source.json

\chapter{Physics problem phyx\_mini\_0286}
\label{ch:PhyXMiniProblems_problem_phyx_mini_0286}

\paragraph{Problem source.}
\#\# Physical scenario

A sand scorpion can detect the motion of a nearby beetle (its prey) by the waves the motion sends along the sand surface. The waves are of two types: transverse waves traveling at \textbackslash{}(v\_t = 50\textbackslash{},\textbackslash{}mathrm\{m/s\}\textbackslash{}) and longitudinal waves traveling at \textbackslash{}(v\_l = 150\textbackslash{},\textbackslash{}mathrm\{m/s\}\textbackslash{}). If a sudden motion sends out such waves, a scorpion can tell the distance of the beetle from the difference \textbackslash{}(\textbackslash{}Delta t\textbackslash{}) in the arrival times of the waves at its leg nearest the beetle.

\#\# Question

If \textbackslash{}(\textbackslash{}Delta t = 4.0\textbackslash{},\textbackslash{}mathrm\{ms\}\textbackslash{}), what is the beetle's distance?

\#\# Answer choices

- A: A: 20 cm

- B: B: 25 cm

- C: C: 35 cm

- D: D: 30 cm

\#\# Figure caption (auxiliary; use the image as primary evidence)

The image depicts a schematic illustration showing two objects, a scorpion and a beetle, positioned on curved surfaces. 

- **Top Section**: A scorpion is standing on an arched surface. Two arrows labeled \textbackslash{}(\textbackslash{}vec\{v\}\_l\textbackslash{}) extend outwards and upwards from the curve, showing movement or propagation direction.

- **Middle Section**: A vertical line marked with \textbackslash{}(d\textbackslash{}) connects the scorpion at the top with the beetle at the bottom, indicating a distance between them.

- **Bottom Section**: A beetle is shown on another arched surface. Two arrows labeled \textbackslash{}(\textbackslash{}vec\{v\}\_t\textbackslash{}) extend outwards and upwards from this curve, similar to the ones near the scorpion.

- **Text**: The bottom part of the illustration includes the label "Beetle," showing the identification of the bottom figure. There are directional vectors, \textbackslash{}(\textbackslash{}vec\{v\}\_l\textbackslash{}) and \textbackslash{}(\textbackslash{}vec\{v\}\_t\textbackslash{}), which could represent forces, velocities, or signals.

Overall, the illustration seems to depict a comparative concept involving the scorpion and beetle.

\#\# Dataset metadata

Wave Properties; Multi-Formula Reasoning, Physical Model Grounding Reasoning

\paragraph{Recorded answer/context.}
D: D: 30 cm

\paragraph{Figure/image path.}
/root/proposal\_for\_physic/science-mango/phyx\_mini\_run/phyx\_data/test\_image/286.png

\paragraph{Formalization target.}
create a compiling Lean file with sorry bodies at `PhyXMiniProblems/problem\_phyx\_mini\_0286.lean`. The Lean declarations must preserve the physical quantities, dimensions or dimensional roles, figure labels, governing-law hypotheses, and final relation expressed by this problem.
Use Mathlib/Physlib names found through LeanExplore where available. If a physics API is missing, introduce faithful local abstractions rather than scalar placeholder aliases.

\begin{theorem}[Physics formalization target]
\label{thm:physics:phyx_mini_0286:target}
\lean{PhyXMiniProblems.ProblemPhyXMini0286.beetleDistance_is_thirtyCentimeters_and_choiceD}
\uses{def:physics:phyx-mini-0286:phyxminiproblems-problemphyxmini0286-lengthinmeters, def:physics:phyx-mini-0286:phyxminiproblems-problemphyxmini0286-lengthincentimeters, def:physics:phyx-mini-0286:phyxminiproblems-problemphyxmini0286-sandwavedetectionsetup, def:physics:phyx-mini-0286:phyxminiproblems-problemphyxmini0286-matchesproblemandfigure, def:physics:phyx-mini-0286:phyxminiproblems-problemphyxmini0286-hasphysicalparameters, def:physics:phyx-mini-0286:phyxminiproblems-problemphyxmini0286-satisfiescommonpathwavetravelmodel, def:physics:phyx-mini-0286:phyxminiproblems-problemphyxmini0286-matchesanswerchoice}
The assigned autoformalize agent should translate this physics problem into Lean declarations in the covered file, with theorem and lemma proof bodies written as `by sorry`.
\end{theorem}
\begin{proof}
This is an autoformalization task, not a proof task. Produce statements that can later be proved by the physics prover without weakening the physical model.
\end{proof}

% --- Archon declaration topology begin ---
\section{Declaration topology}
\begin{definition}[Length Quantity]
  \label{def:physics:phyx-mini-0286:phyxminiproblems-problemphyxmini0286-lengthquantity}
  \lean{PhyXMiniProblems.ProblemPhyXMini0286.LengthQuantity}
  A nonnegative physical separation, independent of the selected length unit
\end{definition}
\begin{proof}
  This is the declared model definition of Length Quantity.
\end{proof}
\begin{definition}[Duration Quantity]
  \label{def:physics:phyx-mini-0286:phyxminiproblems-problemphyxmini0286-durationquantity}
  \lean{PhyXMiniProblems.ProblemPhyXMini0286.DurationQuantity}
  A nonnegative physical duration, independent of the selected time unit
\end{definition}
\begin{proof}
  This is the declared model definition of Duration Quantity.
\end{proof}
\begin{definition}[Speed Quantity]
  \label{def:physics:phyx-mini-0286:phyxminiproblems-problemphyxmini0286-speedquantity}
  \lean{PhyXMiniProblems.ProblemPhyXMini0286.SpeedQuantity}
  A nonnegative physical propagation speed with dimension length per time
\end{definition}
\begin{proof}
  This is the declared model definition of Speed Quantity.
\end{proof}
\begin{definition}[length Readout]
  \label{def:physics:phyx-mini-0286:phyxminiproblems-problemphyxmini0286-lengthreadout}
  \lean{PhyXMiniProblems.ProblemPhyXMini0286.lengthReadout}
  \uses{def:physics:phyx-mini-0286:phyxminiproblems-problemphyxmini0286-lengthquantity}
  Read a physical length in the selected length unit
\end{definition}
\begin{proof}
  This is the declared model definition of length Readout.
\end{proof}
\begin{definition}[duration Readout]
  \label{def:physics:phyx-mini-0286:phyxminiproblems-problemphyxmini0286-durationreadout}
  \lean{PhyXMiniProblems.ProblemPhyXMini0286.durationReadout}
  \uses{def:physics:phyx-mini-0286:phyxminiproblems-problemphyxmini0286-durationquantity}
  Read a physical duration in the selected time unit
\end{definition}
\begin{proof}
  This is the declared model definition of duration Readout.
\end{proof}
\begin{definition}[speed Readout]
  \label{def:physics:phyx-mini-0286:phyxminiproblems-problemphyxmini0286-speedreadout}
  \lean{PhyXMiniProblems.ProblemPhyXMini0286.speedReadout}
  \uses{def:physics:phyx-mini-0286:phyxminiproblems-problemphyxmini0286-speedquantity}
  Read a physical speed in coherent units of length per time
\end{definition}
\begin{proof}
  This is the declared model definition of speed Readout.
\end{proof}
\begin{definition}[length In Meters]
  \label{def:physics:phyx-mini-0286:phyxminiproblems-problemphyxmini0286-lengthinmeters}
  \lean{PhyXMiniProblems.ProblemPhyXMini0286.lengthInMeters}
  \uses{def:physics:phyx-mini-0286:phyxminiproblems-problemphyxmini0286-lengthquantity, def:physics:phyx-mini-0286:phyxminiproblems-problemphyxmini0286-lengthreadout}
  Metre readout of the beetle-to-scorpion separation
\end{definition}
\begin{proof}
  This is the declared model definition of length In Meters.
\end{proof}
\begin{definition}[length In Centimeters]
  \label{def:physics:phyx-mini-0286:phyxminiproblems-problemphyxmini0286-lengthincentimeters}
  \lean{PhyXMiniProblems.ProblemPhyXMini0286.lengthInCentimeters}
  \uses{def:physics:phyx-mini-0286:phyxminiproblems-problemphyxmini0286-lengthquantity, def:physics:phyx-mini-0286:phyxminiproblems-problemphyxmini0286-lengthreadout}
  Centimetre readout used by the four displayed answers
\end{definition}
\begin{proof}
  This is the declared model definition of length In Centimeters.
\end{proof}
\begin{definition}[duration In Seconds]
  \label{def:physics:phyx-mini-0286:phyxminiproblems-problemphyxmini0286-durationinseconds}
  \lean{PhyXMiniProblems.ProblemPhyXMini0286.durationInSeconds}
  \uses{def:physics:phyx-mini-0286:phyxminiproblems-problemphyxmini0286-durationquantity, def:physics:phyx-mini-0286:phyxminiproblems-problemphyxmini0286-durationreadout}
  Second readout used in the constant-speed propagation law
\end{definition}
\begin{proof}
  This is the declared model definition of duration In Seconds.
\end{proof}
\begin{definition}[duration In Milliseconds]
  \label{def:physics:phyx-mini-0286:phyxminiproblems-problemphyxmini0286-durationinmilliseconds}
  \lean{PhyXMiniProblems.ProblemPhyXMini0286.durationInMilliseconds}
  \uses{def:physics:phyx-mini-0286:phyxminiproblems-problemphyxmini0286-durationquantity, def:physics:phyx-mini-0286:phyxminiproblems-problemphyxmini0286-durationreadout}
  Millisecond readout used for the stated arrival-time difference
\end{definition}
\begin{proof}
  This is the declared model definition of duration In Milliseconds.
\end{proof}
\begin{definition}[speed In Meters Per Second]
  \label{def:physics:phyx-mini-0286:phyxminiproblems-problemphyxmini0286-speedinmeterspersecond}
  \lean{PhyXMiniProblems.ProblemPhyXMini0286.speedInMetersPerSecond}
  \uses{def:physics:phyx-mini-0286:phyxminiproblems-problemphyxmini0286-speedquantity, def:physics:phyx-mini-0286:phyxminiproblems-problemphyxmini0286-speedreadout}
  Metre-per-second readout of a sand-wave propagation speed
\end{definition}
\begin{proof}
  This is the declared model definition of speed In Meters Per Second.
\end{proof}
\begin{definition}[Sand Wave Mode]
  \label{def:physics:phyx-mini-0286:phyxminiproblems-problemphyxmini0286-sandwavemode}
  \lean{PhyXMiniProblems.ProblemPhyXMini0286.SandWaveMode}
  The two modes of sand-surface disturbance named in the problem
\end{definition}
\begin{proof}
  This is the declared model definition of Sand Wave Mode.
\end{proof}
\begin{definition}[Detection Endpoint]
  \label{def:physics:phyx-mini-0286:phyxminiproblems-problemphyxmini0286-detectionendpoint}
  \lean{PhyXMiniProblems.ProblemPhyXMini0286.DetectionEndpoint}
  The emitting and detecting endpoints of the timing experiment
\end{definition}
\begin{proof}
  This is the declared model definition of Detection Endpoint.
\end{proof}
\begin{definition}[Figure Object]
  \label{def:physics:phyx-mini-0286:phyxminiproblems-problemphyxmini0286-figureobject}
  \lean{PhyXMiniProblems.ProblemPhyXMini0286.FigureObject}
  The two animals shown in the supplied image
\end{definition}
\begin{proof}
  This is the declared model definition of Figure Object.
\end{proof}
\begin{definition}[Figure Level]
  \label{def:physics:phyx-mini-0286:phyxminiproblems-problemphyxmini0286-figurelevel}
  \lean{PhyXMiniProblems.ProblemPhyXMini0286.FigureLevel}
  Vertical placement of an object or label in the source image
\end{definition}
\begin{proof}
  This is the declared model definition of Figure Level.
\end{proof}
\begin{definition}[Figure Distance Label]
  \label{def:physics:phyx-mini-0286:phyxminiproblems-problemphyxmini0286-figuredistancelabel}
  \lean{PhyXMiniProblems.ProblemPhyXMini0286.FigureDistanceLabel}
  The distance label printed on the central double-headed arrow
\end{definition}
\begin{proof}
  This is the declared model definition of Figure Distance Label.
\end{proof}
\begin{definition}[Figure Velocity Label]
  \label{def:physics:phyx-mini-0286:phyxminiproblems-problemphyxmini0286-figurevelocitylabel}
  \lean{PhyXMiniProblems.ProblemPhyXMini0286.FigureVelocityLabel}
  The velocity-vector labels printed beside the two wave modes
\end{definition}
\begin{proof}
  This is the declared model definition of Figure Velocity Label.
\end{proof}
\begin{definition}[Sand Scorpion Figure]
  \label{def:physics:phyx-mini-0286:phyxminiproblems-problemphyxmini0286-sandscorpionfigure}
  \lean{PhyXMiniProblems.ProblemPhyXMini0286.SandScorpionFigure}
  \uses{def:physics:phyx-mini-0286:phyxminiproblems-problemphyxmini0286-sandwavemode, def:physics:phyx-mini-0286:phyxminiproblems-problemphyxmini0286-figureobject, def:physics:phyx-mini-0286:phyxminiproblems-problemphyxmini0286-figurelevel, def:physics:phyx-mini-0286:phyxminiproblems-problemphyxmini0286-figuredistancelabel, def:physics:phyx-mini-0286:phyxminiproblems-problemphyxmini0286-figurevelocitylabel}
  Qualitative geometry and labels read from the primary bitmap. The locations of the velocity labels are kept as visual evidence only; they do not assert that either wave begins at a different physical source
\end{definition}
\begin{proof}
  This is the declared model definition of Sand Scorpion Figure.
\end{proof}
\begin{definition}[Sand Wave Detection Setup]
  \label{def:physics:phyx-mini-0286:phyxminiproblems-problemphyxmini0286-sandwavedetectionsetup}
  \lean{PhyXMiniProblems.ProblemPhyXMini0286.SandWaveDetectionSetup}
  \uses{def:physics:phyx-mini-0286:phyxminiproblems-problemphyxmini0286-lengthquantity, def:physics:phyx-mini-0286:phyxminiproblems-problemphyxmini0286-durationquantity, def:physics:phyx-mini-0286:phyxminiproblems-problemphyxmini0286-speedquantity, def:physics:phyx-mini-0286:phyxminiproblems-problemphyxmini0286-sandwavemode, def:physics:phyx-mini-0286:phyxminiproblems-problemphyxmini0286-detectionendpoint, def:physics:phyx-mini-0286:phyxminiproblems-problemphyxmini0286-sandscorpionfigure}
  The physical observables in the detection experiment. In particular, beetleDistance is an independent unknown: it is not defined using the recorded answer or either wave speed
\end{definition}
\begin{proof}
  This is the declared model definition of Sand Wave Detection Setup.
\end{proof}
\begin{definition}[Matches Problem And Figure]
  \label{def:physics:phyx-mini-0286:phyxminiproblems-problemphyxmini0286-matchesproblemandfigure}
  \lean{PhyXMiniProblems.ProblemPhyXMini0286.MatchesProblemAndFigure}
  \uses{def:physics:phyx-mini-0286:phyxminiproblems-problemphyxmini0286-durationinmilliseconds, def:physics:phyx-mini-0286:phyxminiproblems-problemphyxmini0286-speedinmeterspersecond, def:physics:phyx-mini-0286:phyxminiproblems-problemphyxmini0286-sandwavedetectionsetup}
  Numerical readouts from the prose and qualitative labels from the primary image. This interface contains no numerical value for beetleDistance
\end{definition}
\begin{proof}
  This is the declared model definition of Matches Problem And Figure.
\end{proof}
\begin{definition}[Has Physical Parameters]
  \label{def:physics:phyx-mini-0286:phyxminiproblems-problemphyxmini0286-hasphysicalparameters}
  \lean{PhyXMiniProblems.ProblemPhyXMini0286.HasPhysicalParameters}
  \uses{def:physics:phyx-mini-0286:phyxminiproblems-problemphyxmini0286-lengthinmeters, def:physics:phyx-mini-0286:phyxminiproblems-problemphyxmini0286-durationinseconds, def:physics:phyx-mini-0286:phyxminiproblems-problemphyxmini0286-speedinmeterspersecond, def:physics:phyx-mini-0286:phyxminiproblems-problemphyxmini0286-sandwavedetectionsetup}
  Positivity and temporal ordering for a nondegenerate detection event
\end{definition}
\begin{proof}
  This is the declared model definition of Has Physical Parameters.
\end{proof}
\begin{definition}[Satisfies Common Path Wave Travel Model]
  \label{def:physics:phyx-mini-0286:phyxminiproblems-problemphyxmini0286-satisfiescommonpathwavetravelmodel}
  \lean{PhyXMiniProblems.ProblemPhyXMini0286.SatisfiesCommonPathWaveTravelModel}
  \uses{def:physics:phyx-mini-0286:phyxminiproblems-problemphyxmini0286-lengthinmeters, def:physics:phyx-mini-0286:phyxminiproblems-problemphyxmini0286-durationinseconds, def:physics:phyx-mini-0286:phyxminiproblems-problemphyxmini0286-speedinmeterspersecond, def:physics:phyx-mini-0286:phyxminiproblems-problemphyxmini0286-sandwavedetectionsetup}
  The ideal common-path time-of-flight model. Each mode obeys distance = speed × travel duration, and simultaneous emission makes the measured lag the transverse duration minus the longitudinal duration. These are governing laws and contain no requested distance value or answer choice
\end{definition}
\begin{proof}
  This is the declared model definition of Satisfies Common Path Wave Travel Model.
\end{proof}
\begin{lemma}[distance In Meters eq arrival Difference div reciprocal Speed Gap]
  \label{lem:physics:phyx-mini-0286:phyxminiproblems-problemphyxmini0286-distanceinmeters-eq-arrivaldifference-div-reciprocalspeedgap}
  \lean{PhyXMiniProblems.ProblemPhyXMini0286.distanceInMeters_eq_arrivalDifference_div_reciprocalSpeedGap}
  \uses{def:physics:phyx-mini-0286:phyxminiproblems-problemphyxmini0286-lengthinmeters, def:physics:phyx-mini-0286:phyxminiproblems-problemphyxmini0286-durationinseconds, def:physics:phyx-mini-0286:phyxminiproblems-problemphyxmini0286-speedinmeterspersecond, def:physics:phyx-mini-0286:phyxminiproblems-problemphyxmini0286-sandwavedetectionsetup, def:physics:phyx-mini-0286:phyxminiproblems-problemphyxmini0286-hasphysicalparameters, def:physics:phyx-mini-0286:phyxminiproblems-problemphyxmini0286-satisfiescommonpathwavetravelmodel}
  For two modes traversing the same distance, the arrival lag gives d = Δt / (1 / v\_t - 1 / v\_l). This is a derived general relation, not a premise of the final result
\end{lemma}
\begin{proof}
  The conclusion follows from the governing relations and figure readouts named in the statement of distance In Meters eq arrival Difference div reciprocal Speed Gap.
\end{proof}
\begin{definition}[Answer Choice]
  \label{def:physics:phyx-mini-0286:phyxminiproblems-problemphyxmini0286-answerchoice}
  \lean{PhyXMiniProblems.ProblemPhyXMini0286.AnswerChoice}
  Labels of the four answers printed with the problem
\end{definition}
\begin{proof}
  This is the declared model definition of Answer Choice.
\end{proof}
\begin{definition}[centimeters]
  \label{def:physics:phyx-mini-0286:phyxminiproblems-problemphyxmini0286-answerchoice-centimeters}
  \lean{PhyXMiniProblems.ProblemPhyXMini0286.AnswerChoice.centimeters}
  \uses{def:physics:phyx-mini-0286:phyxminiproblems-problemphyxmini0286-answerchoice}
  Distance in centimetres printed beside each answer label
\end{definition}
\begin{proof}
  This is the declared model definition of centimeters.
\end{proof}
\begin{definition}[Matches Answer Choice]
  \label{def:physics:phyx-mini-0286:phyxminiproblems-problemphyxmini0286-matchesanswerchoice}
  \lean{PhyXMiniProblems.ProblemPhyXMini0286.MatchesAnswerChoice}
  \uses{def:physics:phyx-mini-0286:phyxminiproblems-problemphyxmini0286-lengthquantity, def:physics:phyx-mini-0286:phyxminiproblems-problemphyxmini0286-lengthincentimeters, def:physics:phyx-mini-0286:phyxminiproblems-problemphyxmini0286-answerchoice}
  A modeled physical distance agrees exactly with a displayed choice
\end{definition}
\begin{proof}
  This is the declared model definition of Matches Answer Choice.
\end{proof}
% --- Archon declaration topology end ---
% --- Archon physics formalization source end ---
